\chapter{Gonorrhoea}

\section*{Definition and Overview}
Gonorrhoea is a sexually transmitted infection (STI) caused by the bacterium \textit{Neisseria gonorrhoeae}. It affects the urethra, cervix, rectum, throat, and eyes, and is spread through unprotected vaginal, anal, and oral sex. The infection often causes no symptoms, particularly in women, so it can be passed on without either partner knowing. If left untreated, gonorrhoea can lead to serious complications such as pelvic inflammatory disease, infertility, and ectopic pregnancy in women, and epididymitis in men. Gonorrhoea is easily treated with antibiotics, but antibiotic resistance is a growing concern, so testing, treatment, and partner notification are essential to control it.

\section*{Epidemiology}
Gonorrhoea is one of the most common bacterial STIs in Australia and worldwide. More than 30,000 cases are reported in Australia each year, and numbers have been rising steadily. The highest rates are in young people aged 15--29, and rates are particularly high among men who have sex with men and among Aboriginal and Torres Strait Islander people. Worldwide, the World Health Organization estimates more than 80 million new infections each year. Because many infections are asymptomatic, the true number of cases is higher than reported, and untreated infection fuels continued transmission.

\section*{Aetiology and Pathophysiology}
Gonorrhoea is caused by infection with \textit{Neisseria gonorrhoeae}, a bacterium that targets the mucous membranes.

\begin{itemize}
    \item \textbf{Transmission:} through sexual contact with infected genital, anal, or throat secretions
    \item \textbf{Sites of infection:} the urethra and cervix, the rectum, the throat, and the conjunctiva of the eye
    \item \textbf{Invasion:} the bacteria attach to mucosal cells, multiply, and trigger an inflammatory response that produces pus and discharge
    \item \textbf{Incubation:} symptoms, when they occur, appear 2--14 days after exposure
    \item \textbf{Asymptomatic carriage:} particularly common in women and with throat and rectal infections
    \item \textbf{Spread within the body:} in women, the infection can ascend to the fallopian tubes causing pelvic inflammatory disease; rarely, it spreads through the bloodstream (disseminated infection)
    \item \textbf{Vertical transmission:} during childbirth, an infected mother can pass gonorrhoea to her baby's eyes
\end{itemize}

\section*{Clinical Features}
The presentation depends on the site of infection and the person's sex.

\begin{itemize}
    \item Men: a thick urethral discharge, burning or pain on urination, and sometimes painful or swollen testicles
    \item Women: often no symptoms; when present, vaginal discharge, pain on urination, pelvic pain, and bleeding between periods or after sex
    \item Rectal infection: discharge, itching, soreness, bleeding, or pain on defecation, or no symptoms at all
    \item Throat infection: usually no symptoms, occasionally a sore throat
    \item Eye infection: in newborns, red, swollen, discharging eyes within a few days of birth; in adults, conjunctivitis after contact with infected secretions
    \item Disseminated infection: fever, skin lesions, and painful swollen joints (rare)
    \item Complications of pelvic inflammatory disease: pelvic pain, fever, and pain during sex
\end{itemize}

\section*{Differential Diagnosis}
The symptoms of gonorrhoea overlap with other STIs and genitourinary conditions.

\begin{itemize}
    \item Chlamydia, which commonly coexists with gonorrhoea
    \item Trichomoniasis, caused by a parasite
    \item Mycoplasma genitalium and ureaplasma infections
    \item Urinary tract infection and cystitis
    \item Bacterial vaginosis and candidiasis (thrush) in women
    \item Non-specific urethritis in men
    \item Appendicitis or ectopic pregnancy in women with pelvic pain
    \item Reactive arthritis (Reiter syndrome)
\end{itemize}

\section*{When to Seek Medical Care}
Get tested if you have symptoms, if a partner has been diagnosed with an STI, or if you have had unprotected sex with a new or casual partner. Regular STI screening is recommended for sexually active people, especially those under 30 and men who have sex with men.

\textbf{Call 000 (or local emergency services) for:}
\begin{itemize}
    \item Severe pelvic or abdominal pain with fever in a woman, which may indicate pelvic inflammatory disease, a tubo-ovarian abscess, or ectopic pregnancy
    \item Sudden severe joint pain and swelling with fever, which may indicate disseminated infection
    \item A baby with red, discharging eyes needs urgent medical review
    \item Signs of sepsis, including high fever, confusion, or collapse
\end{itemize}

\section*{Diagnosis and Investigations}
Testing is simple and accurate, and is usually combined with testing for other STIs.

\begin{itemize}
    \item \textbf{NAAT (nucleic acid amplification test):} the standard test, detecting gonorrhoea DNA in urine or swabs from the urethra, cervix, vagina, rectum, or throat
    \item \textbf{Culture:} used to confirm infection and to test antibiotic sensitivity, particularly for treatment failure or suspected resistance
    \item \textbf{Multiple-site testing:} swabs are taken from sites relevant to the person's sexual practices
    \item \textbf{Co-testing:} testing for chlamydia, syphilis, and HIV is recommended because they often coexist
    \item \textbf{Partner notification:} sexual partners from the last two weeks are advised to be tested and treated
    \item \textbf{Repeat testing:} a test at three months after treatment checks for reinfection
\end{itemize}

\section*{Management}
Gonorrhoea is treated with antibiotics, and cure rates are high when treatment is completed.

\begin{itemize}
    \item \textbf{Antibiotics:} a single injection of ceftriaxone is the standard treatment; oral azithromycin may be added in some guidelines
    \item \textbf{Treatment of partners:} sexual contacts are treated to prevent reinfection and further spread
    \item \textbf{Sexual abstinence:} avoiding sex for at least seven days after treatment and until symptoms resolve
    \item \textbf{Test of cure:} arranged in some cases, particularly for throat infection or when symptoms persist
    \item \textbf{Follow-up testing:} a repeat test at three months is recommended because reinfection is common
    \item \textbf{Pregnancy:} gonorrhoea is treated in pregnancy with safe antibiotics, and the baby's eyes are treated at birth with antibiotic ointment as standard care
    \item \textbf{Management of resistance:} infections that do not respond are reported, and specialist advice is sought
\end{itemize}

\section*{Complications}
\begin{itemize}
    \item In women: pelvic inflammatory disease, chronic pelvic pain, tubal scarring, infertility, and increased risk of ectopic pregnancy
    \item In men: epididymitis and, rarely, reduced fertility
    \item Neonatal eye infection, which can cause blindness if not treated promptly
    \item Disseminated gonococcal infection affecting joints, skin, and rarely the heart or brain
    \item Increased risk of HIV transmission and acquisition
    \item Antibiotic-resistant infection
    \item Psychological distress and effects on relationships
\end{itemize}

\section*{Prognosis and Outcomes}
With early diagnosis and complete treatment, gonorrhoea is cured in nearly all cases, and complications are prevented. The main threats are silent infection causing long-term damage before diagnosis, reinfection from untreated partners, and rising antibiotic resistance. Regular testing, prompt treatment, and partner notification have reduced the burden of gonorrhoea in communities where they are applied consistently. Treated early, gonorrhoea has no lasting effects on fertility or long-term health.

\section*{Prevention and Self-Care}
\begin{itemize}
    \item Use condoms consistently for vaginal, anal, and oral sex
    \item Get tested regularly if you are sexually active, particularly if you are under 30 or have new or multiple partners
    \item Seek testing after any episode of unprotected sex with a new or casual partner
    \item Complete the full course of treatment and attend follow-up testing
    \item Inform sexual partners so they can be tested and treated
    \item Avoid sex until you and your partners have completed treatment
    \item Consider STI testing as a routine part of starting a new relationship
    \item Access sexual health clinics, which provide confidential, free or low-cost testing
\end{itemize}

\section*{Sources and Further Reading}
The following reputable sources provide further information for healthcare students and patients seeking reliable, current advice about gonorrhoea.

\begin{itemize}
    \item Healthdirect Australia. \textit{Gonorrhoea: symptoms, testing, and treatment.} https://www.healthdirect.gov.au/
    \item Melbourne Sexual Health Centre. \textit{Gonorrhoea.} https://www.mshc.org.au/
    \item NSW Sexual Health Infolink. \textit{Gonorrhoea fact sheet.} https://www.shil.nsw.gov.au/
    \item NHS. \textit{Gonorrhoea.} https://www.nhs.uk/
    \item Centers for Disease Control and Prevention. \textit{Gonorrhea.} https://www.cdc.gov/
    \item World Health Organization. \textit{Multi-drug resistant gonorrhoea.} https://www.who.int/
\end{itemize}
